\documentclass{ctexart}
\usepackage{amsmath,amssymb,amsthm}
\usepackage{graphicx}
\usepackage{booktabs}
\usepackage{algorithm}
\usepackage{algpseudocode}
\usepackage{hyperref}
\usepackage{geometry}
\geometry{margin=2.5cm}

\newtheorem{theorem}{定理}[section]
\newtheorem{lemma}[theorem]{引理}
\newtheorem{definition}[theorem]{定义}
\renewcommand{\proofname}{证明}

\title{中英文 \LaTeX{} 测试文档\\zh/en Test Document}
\author{Author Name}
\date{\today}

\begin{document}
\maketitle

\begin{abstract}
本文用于测试 \LaTeX{} 的中英文混排、公式、图表、算法、定理与参考文献等常用功能。
This document tests common \LaTeX{} features for engineering papers.
\end{abstract}

\section{中文测试}
输入中文测试。自动保存。中英文混排：本文使用 \LaTeX{} 排版，English text is mixed here.

\section{English Test}
This is a test for English. The quick brown fox jumps over the lazy dog.

\section{行内与行间公式}
行内公式：$a^2+b^2=c^2$，$\alpha,\beta,\gamma$，$x_i^{(k)}$。

行间公式：
\begin{equation}
\label{eq:integral}
\int_{1}^{5} x^2 \,dx = \frac{125-1}{3} = \frac{124}{3}.
\end{equation}

\begin{equation}
\label{eq:conv}
\int_{-\infty}^{+\infty} e^{-\tau} U(\tau) \,d\tau.
\end{equation}

多行对齐：
\begin{align}
f(x) &= (x+1)^2 \nonumber\\
     &= x^2 + 2x + 1.
\end{align}

矩阵与分段函数：
\[
A=\begin{bmatrix}
1 & 2\\
3 & 4
\end{bmatrix},\qquad
g(x)=
\begin{cases}
x, & x\ge 0,\\
-x, & x<0.
\end{cases}
\]

\section{图表}
\begin{table}[htbp]
\centering
\caption{实验数据}
\label{tab:data}
\begin{tabular}{lcc}
\toprule
方法 & 准确率 & 时间/s\\
\midrule
方法A & 0.91 & 1.2\\
方法B & 0.95 & 2.4\\
\bottomrule
\end{tabular}
\end{table}

如表~\ref{tab:data} 所示，方法B准确率更高。图~\ref{fig:demo} 为示例图。

\begin{figure}[htbp]
\centering
% \includegraphics[width=0.6\textwidth]{figure.png}
\fbox{\rule{0pt}{3cm}\rule{6cm}{0pt}}
\caption{示例图}
\label{fig:demo}
\end{figure}

\section{算法}
\begin{algorithm}[htbp]
\caption{示例算法}
\label{alg:demo}
\begin{algorithmic}[1]
\Require 输入 $x$
\Ensure 输出 $y$
\State 初始化 $y \gets 0$
\For{$i=1$ to $n$}
    \State $y \gets y + x_i$
\EndFor
\State \Return $y$
\end{algorithmic}
\end{algorithm}

\section{定理与证明}
\begin{definition}
若 $f$ 在区间 $I$ 上连续，则称 $f$ 在 $I$ 上连续。
\end{definition}

\begin{theorem}
\label{thm:main}
设 $f$ 在 $[a,b]$ 上连续，则 $f$ 在 $[a,b]$ 上有界。
\end{theorem}

\begin{proof}
由连续函数的性质可知结论成立。
\end{proof}

由定理~\ref{thm:main} 可得相应推论。

\section{列表与引用}
\begin{itemize}
    \item 无序列表项一；
    \item 无序列表项二。
\end{itemize}

\begin{enumerate}
    \item 有序列表项一；
    \item 有序列表项二。
\end{enumerate}

引用示例：见式~\eqref{eq:integral}、表~\ref{tab:data}、图~\ref{fig:demo}、算法~\ref{alg:demo}。

\section{参考文献}
引用示例~\cite{knuth1984}。

\begin{thebibliography}{9}
\bibitem{knuth1984}
D. E. Knuth. \emph{The \TeX book}. Addison-Wesley, 1984.
\end{thebibliography}

\end{document}